\section{Peripheral return phases and a quadratic regularity budget}
\label{sec:peripheral-return-phases}
The spectral phase of a return operator is not a constant phase per
collision. We keep these parameters separate. The physical Fourier
variable is $z\in\R^4$, the return spectral phase is
$\xi\in[-\pi,\pi]$, and the corresponding unit spectral parameter is
$\lambda=e^{i\xi}$. For a section function define
\begin{equation}\label{eq:peripheral-return-defect}
 \mathcal E_{R,z,\xi}f
   =f\circ F_R^*-e^{i(\xi+z\cdot g_R)}f,
       \qquad g_R=G_R-\bar G_R.
\end{equation}
Theorem~\ref{thm:intro-actual-return-defects} concerns $\xi=0$.
We prove a joint small-parameter estimate, including $z=0$, and an
exact lifting identity for every $\xi$. We also improve the regularity
constraint by using the centered second moment when removing the
short collision end blocks.

\subsection{Collision end blocks in the stationary second moment}
\begin{lemma}[Quadratic endpoint budget]
\label{lem:quadratic-endpoint-budget}
There are $a,b,r_0>0$, uniform in $R$, such that
\begin{equation}\label{eq:quadratic-endpoint-budget}
 0<r=|w|\le r_0,\quad H\ge2,\quad r^2(1+\log H)\le a
 \quad\Longrightarrow\quad
 \|\mathcal D_{R,w,\sigma}q\|_1\ge br^2
\end{equation}
for every $\sigma\in\R$ and every collision phase $q$ with
$|q|=1$ and $\|q\|_{\mathrm{BV}}\le H$.
\end{lemma}
\begin{proof}
Use the diffusive middle block of
Lemma~\ref{lem:diffusive-middle-damping}, of length
$m=\lceil A/r^2\rceil$, and the endpoint smoothing in the proof of
Theorem~\ref{thm:near-origin-circle-defect}. Thus
$\epsilon=b_0/H$, $\ell=\lceil b_1+b_2\log H\rceil$, and
$N=m+2\ell$. The lower pairing is at least $1-Nd$, where
$d=\|\mathcal D_{R,w,\sigma}q\|_1$.

The improvement concerns the removal of the phase on the first and
last $\ell$ collisions. Since $h_R$ is centered and its collision
covariances are absolutely summable, uniformly in $R$,
\[
 \|S_{\ell,R}h_R\|_2\le C\sqrt\ell.
\]
All endpoint factors have modulus at most one. Invariance and
$|e^{iu}-e^{iv}|\le|u-v|$ therefore bound the change in the pairing by
\begin{equation}\label{eq:quadratic-endpoint-error}
 |w|\int\left|S_{\ell,R}h_R+
       (S_{\ell,R}h_R)\circ T_R^{m+\ell}\right|\dd\nu
       \le Cr\sqrt\ell.
\end{equation}
No independence of the two end blocks is used. This replaces the
supremum estimate $C r\ell$ in the earlier proof.

After endpoint smoothing and middle-block smoothing, the same
chronological word \eqref{eq:untwisted-end-word} remains. Its two
untwisted powers may be replaced by $\Pi_R$ at cost
$C\epsilon^{-4}\vartheta^\ell$. Its principal pairing has modulus at
most $1/8$. Choose $b_0,b_1,b_2$ exactly as in the earlier proof so
that the two endpoint smoothing and mixing errors total less than
$1/8$. Reduce $r_0$ so that the middle unsmoothing error is less than
$1/16$. Finally choose $a$ so that \eqref{eq:quadratic-endpoint-error}
is less than $1/16$. The original pairing then has modulus at most
$1/2$. The new constraint also gives $N\le C/r^2$, whence
$Nd\ge1/2$ proves \eqref{eq:quadratic-endpoint-budget}.
\end{proof}

\begin{lemma}[Joint small collision frequency and constant phase]
\label{lem:joint-collision-phase-defect}
Let $s=\operatorname{dist}(\sigma,2\pi\Z)$ and
$t=(|w|^2+s^2)^{1/2}$. There are $a,b,t_0>0$ such that
\begin{equation}\label{eq:joint-collision-circle}
 0<t\le t_0,\quad t^2(1+\log H)\le a
 \quad\Longrightarrow\quad
 \|\mathcal D_{R,w,\sigma}q\|_1\ge bt^2
\end{equation}
for every unit-modulus $q$ of variation at most $H\ge2$.
All constants are uniform in $R$.
\end{lemma}
\begin{proof}
Write $r=|w|$ and choose a fixed large number $B$.
If $s\le Br^2$, then $r>0$ and
$t^2\le(1+B^2t_0^2)r^2$. Lemma~\ref{lem:quadratic-endpoint-budget}
proves the claim, after adjusting the constants.

Suppose $s>Br^2$. Centering $q$ in the collision correlation bound
and normalizing its BV norm give
\[
 \left|\int\overline q(q\circ T_R^j)\dd\nu
                   -\left|\int q\dd\nu\right|^2\right|
       \le C H^2e^{-\gamma j}.
\]
Choose $\ell=\lceil b_1+b_2\log H\rceil$ so that this bound is at
most $1/16$ for every $j\ge\ell$. Reduce $t_0$ below $1/4$.
There is an integer
\[
 \ell\le j\le\ell+\lceil2\pi/s\rceil+1
       \quad\hbox{such that}\quad \cos(j\sigma)\le-3/4.
\]
Indeed consecutive phases, in the direction of the representative
of $\sigma$ of absolute value $s$, traverse a full circle in this
many steps; a phase within $s$ of $\pi$ has the required cosine.
This observation includes a negative representative.

The exact iteration of the phase defect implies
\[
 \left|\int\overline q(q\circ T_R^j)\dd\nu
       -e^{ij\sigma}\int e^{iw\cdot S_{j,R}h_R}\dd\nu\right|
       \le j d,\qquad d=\|\mathcal D_{R,w,\sigma}q\|_1.
\]
The centered second moment bounds the difference of the last integral
from one by $C r\sqrt j$. The selected time satisfies
\[
 r^2j\le C(r^2(1+\log H)+r^2/s)\le C(a+B^{-1}).
\]
First choose $B$ large and then $a$ small, so that
$C r\sqrt j\le1/8$. The real part of the first correlation is at
least $-1/16$, whereas that of the second term is at most
$-3/4+1/8$. Hence $jd\ge1/2$.
Finally
\[
 jt^2\le C\bigl(t^2(1+\log H)+r^2/s+s\bigr)
                  \le C(a+B^{-1}+t_0),
\]
so $d\ge bt^2$. This argument also covers $w=0$, $s>0$; thus no
pure constant-phase direction has been lost.
\end{proof}

\begin{proposition}[Joint collision defects for complex functions]
\label{prop:joint-collision-vector-defect}
With $t$ as in Lemma~\ref{lem:joint-collision-phase-defect}, put
$L(H,t)=1+\log(H/t)$. There are uniform $a,b,t_0>0$ such that
\begin{equation}\label{eq:joint-collision-vector}
 0<t\le t_0,\quad t^2L(H,t)\le a,
 \quad \|f\|_2=1,\quad
 \|f\|_\infty+\|f\|_{\mathrm{BV}}\le H
 \quad\Longrightarrow\quad
 \|\mathcal D_{R,w,\sigma}f\|_2\ge\frac{bt^2}{L(H,t)}.
\end{equation}
The function $f$ may vanish.
\end{proposition}
\begin{proof}
Put $e=\|\mathcal D_{R,w,\sigma}f\|_2$ and
$x=\||f|-\int|f|\dd\nu\|_2$. The modulus-variance estimate
\eqref{eq:approximate-modulus-variance} gives
$x^2\le C H^2e^{-\gamma j}+xje$.
For $\eta=\varepsilon_0t^2$, choose $j\le C L(H,t)$ so that the
first term is at most $\eta^2$. Then $x\le\eta+je$.
Lemma~\ref{lem:bv-circle-truncation} provides a unit phase $q$ with
\[
 \|q\|_{\mathrm{BV}}\le C(1+H),\qquad
 \|q-f\|_2\le3\sqrt2(\eta+je).
\]
Its $L^1$ collision defect is at most $e+6\sqrt2(\eta+je)$.
After reducing $a,t_0$, Lemma~\ref{lem:joint-collision-phase-defect}
applies to this phase. Choose $\varepsilon_0$ small enough to absorb
the $\eta$ term. The resulting inequality proves
\eqref{eq:joint-collision-vector}. This is a controlled modulus
argument, not division by an unspecified nonvanishing representative.
\end{proof}

\subsection{The exact spectral phase on the actual tower}
Recall $\eta_R=\one_{Y_R^*}$ and $e_3=(0,0,1,0)$.
For every real $\xi$ define
\begin{equation}\label{eq:peripheral-lift-definition}
 b_{R,\xi}=\eta_R+e^{i\xi}(1-\eta_R),\qquad
 \mathsf L_{R,z,\xi}f=b_{R,\xi}\mathsf L_{R,z}f,
 \qquad \sigma=c_*\xi,\quad w=z-\sigma e_3.
\end{equation}
The multiplier $b_{R,\xi}$ has modulus one and uniformly bounded
variation, independently of $\xi$. In particular its introduction
does not entail a variation estimate for a long return iterate.

\begin{lemma}[Exact peripheral lift and its two normalizations]
\label{lem:exact-peripheral-lift}
For every real $\xi$, every $z\in\R^4$, and $1\le p<\infty$,
whenever the terms are finite,
\begin{align}
 \|\mathsf L_{R,z,\xi}f\|_{L^p(\nu)}^p
     &=c_*\int_{Y_R^*}r_R^*|f|^p\dd\nu_R^*,
                      \label{eq:peripheral-lift-norm}\\
 \|\mathcal D_{R,w,\sigma}(\mathsf L_{R,z,\xi}f)\|_{L^p(\nu)}^p
     &=c_*\|\mathcal E_{R,z,\xi}f\|_{L^p(\nu_R^*)}^p.
                      \label{eq:peripheral-lift-defect}
\end{align}
The defect vanishes at every nontop level of the genuine return tower.
For $Q_{L,\epsilon}$ in \eqref{eq:finite-age-approximation}, the
approximant $b_{R,\xi}Q_{L,\epsilon}$ has the same $L^1$, $L^2$ and
supremum error bounds as in Proposition~\ref{prop:tower-bv-approximation};
its BV bound is at most a fixed multiple of
\eqref{eq:tower-BV-budget}.
\end{lemma}
\begin{proof}
The third collision compensation coordinate is
$h_{R,3}=1-c_*^{-1}\eta_R$. Consequently the real identity
\begin{equation}\label{eq:peripheral-frequency-shift}
 \sigma+w\cdot h_R=z\cdot h_R+\xi\eta_R
\end{equation}
holds, before taking exponentials. On a tower starting at $y\in Y_R^*$,
\[
 \sum_{i=0}^{j-1}\eta_R(T_R^iy)
       =\begin{cases}0,&j=0,\\1,&1\le j\le r_R^*(y).\end{cases}
\]
Thus the lift in \eqref{eq:peripheral-lift-definition} accumulates
exactly the phase in \eqref{eq:peripheral-frequency-shift}. At a
nontop level its collision equation is exact. At the last level its
defect is
\[
 f(F_R^*y)-e^{i(\xi+z\cdot g_R(y))}f(y).
\]
This includes $r_R^*(y)=1$: the initial level is then itself the top.
The disjoint-level integration formula in
Lemma~\ref{lem:exact-tower-defect} gives both identities. The norm
has the height weight; the top defect does not.

Multiplication by $b_{R,\xi}$ preserves pointwise absolute errors.
The BV product inequality gives
$\|b_{R,\xi}Q_{L,\epsilon}\|_{\mathrm{BV}}
 \le C\|Q_{L,\epsilon}\|_{\mathrm{BV}}+CM_f$,
which proves the last assertion. No BV regularity is assumed for the
full exact lift. The construction is periodic in $\xi$, although the
real coordinates $(w,\sigma)$ used to express it change with its lift
from the circle.
\end{proof}

\subsection{A neighborhood of the induced peripheral parameter}
For $H\ge2$, $0<\rho<1/2$, write
\begin{equation}\label{eq:joint-return-budget}
 \Lambda(H,\rho)=1+\log(H/\rho),\qquad
 K(H,\rho)=\Lambda(H,\rho)\log(2+\Lambda(H,\rho)).
\end{equation}

\begin{theorem}[Joint physical and return spectral phases]
\label{thm:joint-return-spectral-defect}
There are uniform constants $a,b,\rho_0>0$ such that, with
\begin{equation}\label{eq:joint-return-domain}
 \xi\in[-\pi,\pi],\qquad
 \rho=(|z|^2+\xi^2)^{1/2},\qquad
 0<\rho\le\rho_0,\qquad \rho^2K(H,\rho)\le a,
\end{equation}
the following estimates hold:
\begin{align}
 |q|=1\text{ on }Y_R^*,\quad \|q^0\|_{\mathrm{BV}}\le H
 &\quad\Longrightarrow\quad
 \|\mathcal E_{R,z,\xi}q\|_{L^1(\nu_R^*)}\ge b\rho^2,
                             \label{eq:joint-return-circle}\\
 \|f\|_{L^2(\nu_R^*)}=1,\quad
 M_f+V_f\le H
 &\quad\Longrightarrow\quad
 \|\mathcal E_{R,z,\xi}f\|_{L^2(\nu_R^*)}
                                  \ge\frac{b\rho^2}{K(H,\rho)}.
                             \label{eq:joint-return-vector}
\end{align}
Zeros of $f$ are allowed. The case $\xi=0$ improves the sufficient
budget of Theorem~\ref{thm:intro-actual-return-defects} from
$|z|K(H,|z|)\le a$ to $|z|^2K(H,|z|)\le a$.
\end{theorem}
\begin{proof}
Put $(w,\sigma)=(z-c_*\xi e_3,c_*\xi)$. For small $\rho$,
$\operatorname{dist}(\sigma,2\pi\Z)=|\sigma|$. The linear map
$(z,\xi)\mapsto(w,\sigma)$ is invertible, with fixed coefficients.
Therefore there are $c_1,C_1>0$ such that
\begin{equation}\label{eq:joint-frequency-comparability}
 c_1\rho\le t=(|w|^2+\sigma^2)^{1/2}\le C_1\rho.
\end{equation}
In particular the line $w=0$ causes no loss: its nonzero constant
collision phase is covered by
Lemma~\ref{lem:joint-collision-phase-defect}.

For a unit section phase, let $Q=\mathsf L_{R,z,\xi}q$.
It has modulus one. Choose accuracy $\eta=\varepsilon_0\rho^2$.
The approximation in Lemma~\ref{lem:exact-peripheral-lift} and
circle truncation give a phase $\widetilde q$ with
\[
 \|\widetilde q-Q\|_1\le4\eta,\qquad
 1+\log(2+\|\widetilde q\|_{\mathrm{BV}})
                           \le C_{\varepsilon_0}K(H,\rho).
\]
Explicitly one may take $L\le C(1+|\log\eta|)$ and
$\epsilon=\eta/(CLH)$. Apply
Lemma~\ref{lem:joint-collision-phase-defect} at $(w,\sigma)$ and
use \eqref{eq:peripheral-lift-defect}. Changing a function by
$4\eta$ changes its $L^1$ collision defect by at most $8\eta$.
Choose $\varepsilon_0$ to absorb this error relative to $b_0c_1^2\rho^2$,
then decrease $a,\rho_0$ so that the required quadratic budget holds.
This proves \eqref{eq:joint-return-circle}.

For a normalized complex section function let
$Q=\mathsf L_{R,z,\xi}f$ and $S=\|Q\|_2$. The exact norm identity
gives $\sqrt{c_*}\le S\le H$. If
$e=\|\mathcal E_{R,z,\xi}f\|_2$, then
\[
 \|\mathcal D_{R,w,\sigma}(Q/S)\|_2
                        =\sqrt{c_*}\,e/S\le e.
\]
Choose now $\eta=\varepsilon_0\rho^2/K(H,\rho)$.
Take $L\le C(1+\log(H/\eta))$ and
$\epsilon=\eta^2/(CH^2L)$ in the $L^2$ approximation and normalize
the resulting approximant. It gives a collision function $v$ with
\begin{align*}
 \|v\|_2&=1,\qquad \|v-Q/S\|_2\le2\eta,\\
 1+\log(2+\|v\|_\infty+\|v\|_{\mathrm{BV}})+|\log t|
       &\le C(1+|\log\varepsilon_0|)^2K(H,\rho).
\end{align*}
These inequalities follow directly by taking logarithms in
\eqref{eq:tower-BV-budget};
$\log K(H,\rho)\le C\Lambda(H,\rho)$ and
\eqref{eq:joint-frequency-comparability} are used here.
They display the dependence on the accuracy choice.
Proposition~\ref{prop:joint-collision-vector-defect} yields
\[
 e+4\eta\ge
 \frac{b_0c_1^2\rho^2}
      {C(1+|\log\varepsilon_0|)^2K(H,\rho)}.
\]
Fix $\varepsilon_0$ so small that the accuracy term is absorbable;
this is possible since
$\varepsilon_0(1+|\log\varepsilon_0|)^2\to0$.
Only after that choice reduce $a,\rho_0$ to satisfy the collision
quadratic budget. This proves \eqref{eq:joint-return-vector} with
no circular choice of a regularity constant.
\end{proof}

\begin{corollary}[Peripheral annuli with enlarged regularity classes]
\label{cor:peripheral-annulus-budgets}
Consider the physical annulus and its rescaling
\begin{equation}\label{eq:peripheral-annulus-scales}
 2n^{-99/200}\le|z|\le n^{-2/5},\qquad
 2n^{1/200}\le|\sqrt n\,z|\le n^{1/10}.
\end{equation}
For $H\le H_0n^\zeta$, there is $a_0>0$ such that the preceding
theorem holds throughout this annulus and
\[
 |\xi|\le\frac{a_0}{\sqrt{\log(2+n)\log\log(e^e+n)}}
\]
for all sufficiently large $n$. Its complex-function lower bound is
at least $b(|z|^2+\xi^2)/[\log(2+n)\log\log(e^e+n)]$.

More generally, fix $\gamma>0$ and
$0<\beta<\min\{4/5,2\gamma\}$. For
$H\le H_0\exp(n^\beta)$ and $|\xi|\le n^{-\gamma}$, the theorem
holds uniformly on \eqref{eq:peripheral-annulus-scales}, with
complex-function lower bound at least
$b(|z|^2+\xi^2)/(n^\beta\log(2+n))$.
\end{corollary}
\begin{proof}
The lower annulus endpoint gives $\rho\ge2n^{-99/200}$.
For polynomial budgets this implies
$K(H,\rho)\le C\log(2+n)\log\log(e^e+n)$.
Choose $a_0$ small enough that the $\xi^2K$ term meets
\eqref{eq:joint-return-domain}; the $|z|^2K$ term tends to zero.
For the exponential budget, $K(H,\rho)\le Cn^\beta\log(2+n)$,
whereas $\rho^2\le n^{-4/5}+n^{-2\gamma}$.
The strict inequalities on $\beta$ make their product tend to zero.
The lower bounds follow from \eqref{eq:joint-return-vector}.
\end{proof}

\paragraph{Spectral scope.}
The lift in Lemma~\ref{lem:exact-peripheral-lift} treats every return
phase exactly. The quantitative conclusion, in contrast, is on the
joint region \eqref{eq:joint-return-domain}; it is not a bound on the
entire unit circle at a growing regularity budget. The pure spectral
direction $z=0$, $\xi\ne0$ and the shifted-frequency cancellation line
are included. Theorem~\ref{thm:joint-return-spectral-defect} remains a
physical function theorem. The next section gives a precise
operator-defect comparison for finite-rank compressions, without
identifying them with an anisotropic realization of the unmodified
return dynamics.
